% !TEX root = 10.tex
\documentclass[a4paper,10pt]{article}

\usepackage{pdfsync}
\usepackage[utf8]{inputenc}
\usepackage{microtype}
\usepackage{mathtools}
\usepackage{amsthm,amsfonts,amsmath,amssymb}
\usepackage{mathabx}
% \usepackage{MnSymbol}
\usepackage{hyperref}
\usepackage{bbold}
\usepackage[textwidth=2cm,textsize=small]{todonotes}
\usepackage{subcaption}
\usepackage{enumerate}
\usepackage{xspace}
\usepackage{stmaryrd}
\usepackage{verbatim}
\usepackage{cleveref}

\hypersetup{
    colorlinks,
    linkcolor={red!50!black},
    citecolor={blue!50!black},
    urlcolor={blue!80!black}
}

\author{}

\newcommand{\ignore}[1]{}

\newcommand{\st}{s.t.}
\newcommand{\ie}{i.e.}
\newcommand{\eg}{e.g.}

\newcommand{\sep}{\;|\;}
\renewcommand{\rule}[2]{\dfrac{#1}{#2}}
\newcommand{\set}[1]{\left\{#1\right\}}
\newcommand{\setof}[2]{\set{#1 \;\middle|\; #2}}
\newcommand{\tuple}[1]{\left( #1 \right)}
\newcommand{\pair}[2]{\tuple {#1, #2}}
\newcommand{\closure}[3]{\tuple {#1, #2, #3}}
\newcommand{\sem}[1]{\llbracket{#1}\rrbracket}
\newcommand{\powerset}[1]{\mathcal P(#1)}
\newcommand{\map}[2]{\mathsf{map}\; #1\; #2}
\newcommand{\filter}[2]{\mathsf{filter}\; #1\; #2}
\newcommand{\floor}[1]{\left\lfloor #1 \right\rfloor}
\newcommand{\PreRel}[2]{\mathsf{Pre}_{#1}(#2)}
\newcommand{\Pre}[1]{\PreRel {} {#1}}
\newcommand{\PreStar}[1]{\mathsf{Pre}^*(#1)}
\newcommand{\Post}[1]{\mathsf{Post}(#1)}
\newcommand{\PostStar}[1]{\mathsf{Post^*}(#1)}
\newcommand{\Conf}{\mathsf{Conf}}
\newcommand{\upwardclosure}[1]{#1\uparrow}
\newcommand{\card}[1]{\left| #1 \right|}
\newcommand{\lang}[1]{L(#1)}
\newcommand{\goesto}[1]{\xrightarrow{#1}}
\newcommand{\e}{\varepsilon}

\newcommand{\N}{\mathbb N}
\newcommand{\Z}{\mathbb Z}
\newcommand{\Q}{\mathbb Q}

% LTL
\newcommand{\true}{\mathbf{true}}
\newcommand{\false}{\mathbf{false}}
\newcommand{\X}{\mathbf X}
\newcommand{\U}{\mathbf U}
\newcommand{\Release}{\mathbf R}
\newcommand{\F}{\mathbf F}
\newcommand{\G}{\mathbf G}

% CTL
\newcommand{\EX}{\mathbf{EX}}
\newcommand{\AX}{\mathbf{AX}}
\newcommand{\EF}{\mathbf{EF}}
\newcommand{\AF}{\mathbf{AF}}
\newcommand{\EG}{\mathbf{EG}}
\newcommand{\AG}{\mathbf{AG}}
\newcommand{\EU}[2]{\mathbf E(#1 \U #2)}
\newcommand{\AU}[2]{\mathbf A(#1 \U #2)}
\newcommand{\ER}[2]{\mathbf E(#1 \Release #2)}
\newcommand{\AR}[2]{\mathbf A(#1 \Release #2)}

\newcommand{\NL}{\textsf{NL}}
\newcommand{\PTIME}{\textsf{PTIME}}
\newcommand{\PSPACE}{\textsf{PSPACE}}

\makeatletter

\def\dasharrowfill@#1#2#3#4{%
        $\m@th
        \thickmuskip0mu
        \medmuskip\thickmuskip
        \thinmuskip\thickmuskip
        \relax
        #4#1\mkern2mu
        \xleaders\hbox{$#4\mkern2mu#2\mkern2mu$}\hfill
        \mkern2mu
        #3$%
}

\def\dashleftarrowfill@{\dasharrowfill@\leftarrow\relbar\relbar}
\def\dashrightarrowfill@{\dasharrowfill@\relbar\relbar\rightarrow}
\def\dashleftrightarrowfill@{\dasharrowfill@\leftarrow\relbar\rightarrow}
\def\dashLeftarrowfill@{\dasharrowfill@\Leftarrow\Relbar\Relbar}
\def\dashRightarrowfill@{\dasharrowfill@\Relbar\Relbar\Rightarrow}
\def\dashLeftrightarrowfill@{\dasharrowfill@\Leftarrow\Relbar\Rightarrow}

\providecommand*\xdashleftarrow[2][]{%
  \ext@arrow 0055{\dashleftarrowfill@}{#1}{#2}}
\providecommand*\xdashrightarrow[2][]{%
  \ext@arrow 0055{\dashrightarrowfill@}{#1}{#2}}
\providecommand*\xdashleftrightarrow[2][]{%
  \ext@arrow 0055{\dashleftrightarrowfill@}{#1}{#2}}
\providecommand*\xdashLeftarrow[2][]{%
  \ext@arrow 0055{\dashLeftarrowfill@}{#1}{#2}}
\providecommand*\xdashRightarrow[2][]{%
  \ext@arrow 0055{\dashRightarrowfill@}{#1}{#2}}
\providecommand*\xdashLeftrightarrow[2][]{%
  \ext@arrow 0055{\dashLeftrightarrowfill@}{#1}{#2}}

\def\slashedarrowfill@#1#2#3#4#5{%
$\m@th\thickmuskip0mu\medmuskip\thickmuskip\thinmuskip\thickmuskip
  \relax#5#1\mkern-7mu%
  \cleaders\hbox{$#5\mkern-2mu#2\mkern-2mu$}\hfill
  \mathclap{#3}\mathclap{#2}%
  \cleaders\hbox{$#5\mkern-2mu#2\mkern-2mu$}\hfill
  \mkern-7mu#4$%
}

\def\rightslashedarrowfill@{%
  \slashedarrowfill@\relbar\relbar{\raisebox{1.2pt}{$\scriptscriptstyle\diagup$}}\rightarrow}
\newcommand\xslashedrightarrow[2][]{%
  \ext@arrow 0055{\rightslashedarrowfill@}{#1}{#2}}
\makeatother

\newtheorem{exercise}{Exercise}

\newenvironment{solution}{\begin{proof}[\protect{\textnormal{\emph{Solution.}}}]}{\end{proof}}

\let\oldqedhere\qedhere

\newif\ifstartedinmathmode
\renewcommand*{\qedhere}{%
  \relax\ifmmode\startedinmathmodetrue\else\startedinmathmodefalse\fi%
  \ifstartedinmathmode\tag*{\protect\oldqedhere}\else\ifinlist{\hfill\oldqedhere}{\oldqedhere}\fi%
}

\crefname{equation}{}{}
\crefname{exercise}{Exercise}{Exercises}


\title{Formal verification - Tutorial 10}
\date{25/05/2026}

\begin{document}

\maketitle

\section*{Rational series}

Let $\Sigma$ be a finite alphabet.
A \emph{series} is a mapping $f : \Sigma^* \to \Q$.
The \emph{support} of a series $f$ is the language $\supp f = \setof{ w \in \Sigma^*} {f(w) \neq 0}$.
Recall that the set of rational series coincides with the set of series recognised by weighted finite automata.

\begin{exercise}
    Show an example of a rational series with a non-regular support.
\end{exercise}
\begin{solution}
    Consider the series mapping a word $w \in \Sigma^*$
    to the difference of the number of occurrences of $a$ and $b$ in $w$.
    The support of this series is the language of words with an unequal number of $a$'s and $b$'s,
    which is well-known not to be regular.
    The series is rational, in fact a weighted automaton of dimension two suffices to recognise it.
\end{solution}

The \emph{left shift} (or \emph{left quotient}, or \emph{left derivative}) of a series $f$ by a word $w \in \Sigma^*$
is the series $w^{-1} f$ defined by $(w^{-1} f)(u) = f(wu)$ for all $u \in \Sigma^*$.
%
The \emph{rank} of a series $f$ is the dimension of the vector space generated by
the set of all left shifts of $f$, i.e., the set $\setof{ w^{-1} f }{ w \in \Sigma^*}$.
%
Recall that the rank of a series is finite if, and only if, the series is rational.

\begin{exercise}
    Assume that $f, g$ are rational series with rank $k$, resp., $\ell$.
    %
    \begin{enumerate}
        \item Show that the rank of $f + g$ is at most $k + \ell$.
        \item The \emph{Hadamard product} of $f$ and $g$
        is the series $f \odot g$ defined by $(f \odot g)(w) = f(w) g(w)$ for all $w \in \Sigma^*$.
        %
        Show that the rank of $f \odot g$ is at most $k \cdot \ell$.
    \end{enumerate}
\end{exercise}
% \begin{solution}
% \end{solution}

A weighted automaton is \emph{observable} at an initial condition $u \in \Q^{1 \times k}$
if, for every initial condition $u' \in \Q^{1 \times k}$,
$\sem{A_u} = \sem {A_{u'}}$ implies $u = u'$.
%
A weighted automaton is \emph{observable} if it is observable at every initial condition.

 \begin{exercise}
    Show that the observability problem for weighted automata is decidable in polynomial time.
 \end{exercise}

A \emph{parametric weighted automaton} is a tuple
%
\begin{align*}
    A = \tuple {\Sigma, k, \ell, u, M, v}
\end{align*}
%
where
%
\begin{itemize}
    \item $\Sigma$ is a finite input alphabet,
    \item $k \in \N$ is the dimension of the state
    \item $\ell \in \N$ is the number of parameters,
    \item $u \in \poly \Q {p_1, \dots, p_\ell}^{1 \times k}$ is the parametric initial vector,
    \item $M : \Sigma \to \poly \Q {p_1, \dots, p_\ell}^{k \times k}$
    is the \emph{transition function} mapping each letter to a parametric transition matrix, and
    \item $v \in \poly \Q {p_1, \dots, p_\ell}^{k \times 1}$ is the parametric final vector.
\end{itemize}
%
For every given valuation for the parameters $\nu \in \Q^\ell$
we obtain a concrete weighted automaton $A(\nu)$ in the expected way.
This concrete automaton will then recognise the rational series $\sem{A(\nu)} : \Sigma^* \to \Q$.

\begin{exercise}
    Show that the set of all parameter valuations $\nu \in \C^\ell$ \st~$A(\nu)$ recognises the zero series is an algebraic variety.
    Show that polynomial generators of degree polynomial in $k$ suffice to define this variety.
\end{exercise}

The \emph{parameter identifiability problem} for a parametric weighted automaton $A$
amounts to deciding whether, for all parameter valuations $\nu, \nu' \in \C^\ell$,
$\sem {A(\nu)} = \sem {A(\nu')}$ implies $\nu = \nu'$.

 \begin{exercise}
    Show that the parameter identifiability problem is decidable for parametric weighted automata.
 \end{exercise}

\bibliographystyle{plainurl}
\bibliography{bibliography}

\end{document}